\documentclass{article} % For LaTeX2e
\usepackage{iclr2026_conference,times}

% Optional math commands from https://github.com/goodfeli/dlbook_notation.
\input{math_commands.tex}

\usepackage[utf8]{inputenc}
\usepackage{amsmath}
\usepackage{amssymb}
\usepackage{booktabs}
\usepackage{graphicx}
\usepackage{hyperref}
\usepackage{url}
\usepackage{tikz}

\title{Is Spike-Driven Self-Attention Necessary?\\ The Inefficiency of Jaccard Attention in Spiking Sentence Embeddings}

\author{Anonymous Authors \\
\textit{Affiliation}\\
\texttt{email@domain.com}
}

\newcommand{\fix}{\marginpar{FIX}}
\newcommand{\new}{\marginpar{NEW}}

\iclrfinalcopy 
\begin{document}

\maketitle

\begin{abstract}
Spiking Neural Networks (SNNs) provide an energy-efficient alternative to conventional Deep Neural Networks by utilizing discrete, event-driven temporal spikes. However, adopting Transformer-based paradigms into the spiking domain typically requires mapping complex Self-Attention mechanisms into spike operations. In this paper, we present an in-depth mathematical formulation and empirical ablation of a Spike-Driven Self-Attention layer—modeled around a Jaccard-like spike intersection metric—to evaluate its true utility in semantic sentence embedding tasks. Through comprehensive knowledge distillation, we demonstrate that incorporating Jaccard Attention yields a severe execution bottleneck and triples the total spike emission, while offering a negligible absolute accuracy gain (+0.0048 Pearson correlation). We propose a purely recurrent, attention-free Spiking Sentence Embedder operating exclusively through Leaky Integrate-and-Fire (LIF) pooling and Backpropagation Through Time (BPTT). Our minimalist architecture achieves a highly competitive Pearson correlation of 0.8031 on the ALL-STS benchmark with an estimated energy savings ratio of 1298$\times$ compared to standard Transformer MACs, conclusively demonstrating that spatial routing via self-attention is redundant for high-fidelity spiking semantic representations.
\end{abstract}

\section{Introduction}

The computational overhead of Dense Sentence Embedders (e.g., SBERT, SimCSE) is heavily dominated by the dense matrix multiplications required in standard Scaled Dot-Product Attention \citep{reimers2019sentence, gao2021simcse}. Spiking Neural Networks (SNNs) alleviate this constraint by replacing power-hungry Multiply-Accumulate (MAC) operations with highly sparse, addition-only Synaptic Operations (SOPs). Recent neuromorphic NLP research has focused on translating Self-Attention to SNNs \citep{li2022spikeformer, zhu2023spikegpt}, such as the Sparse Coincidence-Based Semantic Attention \citep{coincidence2026} which relies on spike overlaps. However, the fundamental question persists: \textit{Does the temporal recurrent dynamics of an SNN naturally obviate the need for explicit $\mathcal{O}(L^2)$ spatial attention mechanisms?}

In this study, we systematically deconstruct the forward dynamics and learning mechanisms of an SNN sentence embedder down to its core computational primitives. We mathematically define a Jaccard-based Spiking Attention algorithm and contrast it against an attention-free Recurrent Pooler. Our empirical findings across a multi-phase evaluation indicate that the temporal integration capabilities of the recurrent pooler, when optimized via Backpropagation Through Time (BPTT) and Surrogate Gradients, are overwhelmingly sufficient for encoding semantic structures.

\section{Methodology: Architectural Dynamics and BPTT}

Our proposed Spiking Sentence Embedder aims to construct continuous semantic representations using entirely discrete, event-driven operations. The pipeline consists of Spike Encoding, a comparative Attention vs. Pooling mechanism, and Knowledge Distillation.

\subsection{Spike Encoding and Heterogeneous Dynamics}
Raw text is tokenized into vocabulary indices. The Spiking Embedding Layer translates these tokens into temporal spike trains $S_{emb}^{(t)}$ over discrete sequence steps $t \in [1, L]$, where the temporal dimension maps perfectly to sequence length. The membrane potential $V_{emb}^{(t)}$ is modeled via Leaky Integrate-and-Fire (LIF) dynamics:
\begin{equation}
V_{emb}^{(t)} = \min\left(1.0, \beta_{emb} V_{emb}^{(t-1)} + X W_{emb}\right) - S_{emb}^{(t)} \theta_{emb}
\end{equation}
When the pre-reset potential exceeds the threshold $\theta_{emb}$, a spike $S_{emb}^{(t)} = 1$ is fired. We utilize a \textit{Soft-Reset} strategy (subtracting the threshold) rather than a hard reset to absolute zero. Preserving these fractional sub-threshold residual potentials prevents catastrophic semantic information loss across time. Furthermore, we introduce biological heterogeneity: the leak decay $\beta_{emb}$ and thresholds $\theta_{emb}$ are initialized with unique, uniformly distributed values across dimensions (e.g., $\beta \in [1 - 2^{-2}, 1 - 2^{-7}]$). This forces multi-timescale temporal processing immediately upon sequence ingestion.

\subsection{Ablation Baseline: Spike-Driven Jaccard Attention}
To evaluate the impact of spatial routing, we formalized a Spike-Driven Self-Attention mechanism as our ablation baseline. Unlike floating-point dot products, this method computes affinities purely through Boolean spike intersections (Jaccard index numerator).
Embedding spikes $S_{emb}$ are projected into Query ($Q$), Key ($K$), and Value ($V$) spikes using independent LIF neurons. For any query token $i$ and key token $j$, the raw overlap is captured via a match count:
\begin{equation}
M_{i,j}^{(t)} = \sum_{d=1}^{D_{model}} \left( Q_{i,d}^{(t)} \wedge K_{j,d}^{(t)} \right)
\end{equation}
These integer match counts are normalized by the maximum active sequence matches $\max_k (M_{i,k}^{(t)} + \epsilon)$. To maintain the purely spiking pipeline, these normalized fractional scores are passed into an auxiliary set of LIF neurons to generate attention mask spikes $S_{scores}$, which finally gate the Value spikes $V$ via accumulation.
While this avoids dense MACs, calculating $\mathcal{O}(L^2 \cdot D_{model})$ Boolean intersections at every single time-step, followed by dual-layered LIF dynamics, imposes severe computational bottlenecks and dramatically inflates internal spike activity.

\subsection{The Superior Recurrent Pooler (Proposed Architecture)}
We propose that spatial attention is functionally redundant. In our purely recurrent architecture, $S_{emb}^{(t)}$ bypasses attention and directly feeds into a Spiking Dense BPTT Pooler layer. This layer utilizes an Add-Only synaptic operation (since input spikes are exactly $\{0,1\}$).
The final continuous sentence embedding $E$ is extracted by accumulating the fractional membrane potentials across the entire temporal sequence $T$:
\begin{equation}
E = \sum_{t=1}^T V_{pool}^{(t)}
\end{equation}

\textbf{Learning via BPTT and Surrogate Gradients:} The true power of this architecture stems from its backpropagation mechanics. Because the Heaviside step function lacks a valid derivative, we employ a Rectangular Surrogate Mask to propagate errors:
\begin{equation}
\sigma'(V) = \begin{cases} 1 & \text{if } |V - \theta| < \text{window\_size} \\ 0 & \text{otherwise} \end{cases}
\end{equation}
During Backpropagation Through Time, the error $\delta^{(t)}$ at sequence step $t$ is explicitly linked to future time steps via the heterogeneous decay $\beta_{pool}$:
\begin{equation}
\tilde{\delta}^{(t)} = \left( \delta_{out}^{(t)} + \delta^{(t+1)} \cdot \beta_{pool} \right) \cdot \sigma'(V_{pool}^{(t)})
\end{equation}
This explicit temporal unrolling, using the Add-Only Delta rule, allows the $\beta$ parameter to act as a natural contextual router. The pooler inherently learns long-term syntactic dependencies simply by preserving or decaying gradients, rendering the spatial $\mathcal{O}(L^2)$ matrix completely unnecessary.

\subsection{Knowledge Distillation}
To optimize these dynamics, we distill dark knowledge from a pre-trained dense Teacher. The objective minimizes the difference between the Cosine Similarity of the SNN representations and the continuous Teacher score via a Contrastive Hebbian trace margin ($m=0.05$):
\begin{equation}
\mathcal{L} = \left| \cos(E_1, E_2) - \text{Score}_{Teacher} \right|
\end{equation}
Distillation stabilizes the discrete parameter space significantly faster than unsupervised sparse objective learning.

\section{Experimental Evaluation}

The models are implemented in Rust to guarantee bit-accurate MAC-to-SOP translation, avoiding PyTorch simulation artifacts. We trained over 20 epochs across 29,720 validation pairs from ALL-STS.

\subsection{Phase 1: Dimensionality Scaling}
Using the Distillation strategy with the Recurrent Pooler, we scaled the embedding dimension $D_{model}$ to analyze representation capacity. 

\begin{table}[h]
\centering
\caption{Phase 1: Dimensionality Scaling (Recurrent Pooler)}
\begin{tabular}{lcccc}
\toprule
$D_{model}$ & \textbf{ALL-STS} & \textbf{Inference} & \textbf{Est. SOPs/} & \textbf{Energy} \\
& \textbf{Pearson} & \textbf{Latency (ms/pair)} & \textbf{Sentence} & \textbf{Savings} \\
\midrule
32 & 0.6713 & 0.718 & 17,284 & $82,991\times$ \\
64 & 0.7091 & 0.911 & 79,505 & $18,042\times$ \\
128 & 0.8038 & 1.341 & 339,181 & $4,229\times$ \\
\textbf{256} & \textbf{0.8079} & \textbf{1.946} & \textbf{1,407,442} & \textbf{$1,019\times$} \\
\bottomrule
\end{tabular}
\end{table}

The SNN scales impeccably to 256 dimensions, shattering the 0.80 Pearson correlation barrier, a massive improvement over prior Spike Coincidence Attention baselines \citep{coincidence2026} ($r=0.758$).

\subsection{Phase 2: Ablation of Spike-Driven Attention}
Fixing the dimension to $D_{model} = 256$, we performed an architectural ablation against the Jaccard Attention and the minimalist Recurrent Pooler.

\begin{table}[h]
\centering
\caption{Phase 2: Ablation of Spike-Driven Attention ($D_{model}=256$)}
\begin{tabular}{lcccc}
\toprule
\textbf{Architecture} & \textbf{ALL-STS} & \textbf{Inference} & \textbf{Avg Spikes/} & \textbf{Energy} \\
\textbf{Variant} & \textbf{Pearson} & \textbf{(ms/pair)} & \textbf{Sent.} & \textbf{Savings} \\
\midrule
\textbf{No-Attention (Recurrent)} & 0.8031 & \textbf{0.932} & \textbf{2,157} & \textbf{1298$\times$} \\
\textbf{Jaccard Attention} & \textbf{0.8079} & 1.946 & 6,737 & 1019$\times$ \\
Homogeneous-Init (Attn) & 0.7996 & 1.924 & 6,687 & 1179$\times$ \\
\bottomrule
\end{tabular}
\end{table}

\textbf{Discussion:} The empirical data definitively undermines the necessity of Spike-Driven Self-Attention. While Jaccard Attention yields a negligible raw accuracy bump ($\Delta = +0.0048$), it emits over $3\times$ the amount of internal spikes and requires double the execution latency ($1.946$ ms vs $0.932$ ms) per pair. The BPTT gradients of the Recurrent Pooler's temporal context ($ \beta V_{pool}^{(t-1)} $) prove highly expressive, cementing recurrent SNNs as the dominant efficient topology for sequences.

\section{Limitations}
While our attention-free recurrent SNN achieves remarkable sparsity, the peak correlation ($0.8031$) remains marginally below the theoretical limit of the full-precision Teacher model ($>0.85$), indicating an inherent compression bottleneck in discrete binary signaling. Furthermore, our empirical sequence length is capped at $L=128$; verifying temporal decay limits over much longer contexts (e.g., document-level embeddings) requires future investigation.

\section{Conclusion}
Through detailed mathematical formulation and rigorous low-level structural ablation, we demonstrate that forcing Spatial Attention onto Spiking Neural Networks produces an unjustifiable computational bottleneck. The biological temporal decay properties inherent to LIF neurons, when trained via Surrogate Gradients and BPTT, naturally encapsulate complex sentence semantics without requiring quadratic sequence routing. We conclude that Spike-Driven Self-Attention is redundant for dense semantic representations, paving the way for ultra-efficient, purely recurrent SNN embedders.

\section{Data and Code Availability}
To ensure reproducibility, the raw Rust SNN logic, distilled datasets, and mathematical formulations are made available. 

\bibliography{iclr2026_conference}
\bibliographystyle{iclr2026_conference}

\end{document}
